\pagestyle{corporativo}

\begin{frontmatter}

	\title{\projecttitle}

	\author[inst1]{J. Arboleda\corref{cor1}}
	\ead{[correo@dmlsas.com]}

	\author[inst1]{F. Navia}

	\author[inst1]{H. Rosero}

	\cortext[cor1]{Autor correspondiente}
	\address[inst1]{DML Ingenieros Consultores S.A.S., Cali, Colombia, 760025}

	\begin{abstract}
		\thispagestyle{corporativo}
		Se presenta el an\'{a}lisis de flexibilidad de la l\'{i}nea de la bomba de
		vac\'{i}o nueva (l\'{i}nea SIM-007, job CAESAR P2603-PR-SIM-007) del
		proyecto P2603 SW-K60 para Smurfit Westrock, ejecutado conforme a
		ASME B31.3-2016 con el software CAESAR II 2019. La l\'{i}nea, construida en
		acero inoxidable ASTM A312 TP 304L Sch 10S en di\'{a}metros de 4 a 8~NPS,
		opera con agua a 1241~\si{kPa}~g y 90~\si{\celsius}. El modelo se
		gener\'{o} desde el isom\'{e}trico de AutoCAD Plant 3D 2027 mediante
		archivo PCF corregido, y se evaluaron nueve casos de carga (operaci\'{o}n,
		sostenidos y expansi\'{o}n). El esfuerzo m\'{a}ximo de c\'{o}digo fue
		\num{191915.1}~\si{kPa} frente a un admisible de
		\num{273789.8}~\si{kPa}, correspondiente al nodo 100 en el caso de
		expansi\'{o}n 9 (EXP) L9=L1-L3, con una relaci\'{o}n de utilizaci\'{o}n de
		70.1~\%. La evaluaci\'{o}n de cumplimiento de c\'{o}digo arroj\'{o}
		\textbf{CODE COMPLIANCE EVALUATION PASSED}: la l\'{i}nea es aceptable
		seg\'{u}n ASME B31.3-2016 y es la m\'{a}s exigida del paquete de
		flexibilidad del proyecto, por lo que se recomienda conservar la
		configuraci\'{o}n analizada.
	\end{abstract}

	\begin{keyword}
		flexibilidad de tuber\'{i}as \sep CAESAR II \sep ASME B31.3 \sep
		esfuerzo de expansi\'{o}n \sep desplazamientos \sep an\'{a}lisis de esfuerzos
	\end{keyword}

\end{frontmatter}

\pagestyle{corporativo}
\thispagestyle{corporativo}
